% Video voiceover script for the 13-slide presentation
% (Project_Video_Presentation.pptx). One block per slide, written to be
% read aloud verbatim. Figures are given as numerals so they are quick to
% read off the page. Compile with pdfLaTeX on Overleaf.
\documentclass[a4paper,11pt]{article}

\usepackage[a4paper,top=2.4cm,bottom=2.4cm,left=2.4cm,right=2.4cm]{geometry}
\usepackage{newtxtext,newtxmath}
\usepackage{enumitem}
\usepackage{xurl}
\usepackage[hidelinks]{hyperref}

\newcommand{\slideblock}[3]{%
  \section*{Slide #1: #2 \hfill \normalfont\small(\textit{#3})}}

\title{\bfseries Video Voiceover Script\\
\large Does Machine Learning Beat the Credit Scorecard? (Trimester 8 Term Project)}
\author{François Schmitt (Roll Number: 23035010523)\\
\ttfamily f.schmitt@op.iitg.ac.in}
\date{}

\begin{document}
\maketitle

\noindent\textit{Reading guide: each block below is read while the corresponding slide is on screen. Times in the headers are targets; the total is about twelve minutes. Keep the face visible throughout, per the submission guide. Words in [brackets] are stage directions, not to be read.}

\slideblock{1}{Title}{about 35 seconds}
Hello. My name is François Schmitt, I am a student of the B.Sc.\ Honours Data Science and Artificial Intelligence Online Degree Programme at IIT Guwahati, and this is my Trimester 8 term project. The title is: Does Machine Learning Beat the Credit Scorecard? It is a study in credit risk, on real loan data, and it asks whether the machine-learning advantage that research papers report actually survives the criteria that a real lender cares about. Let me show you what I did and what I found.

\slideblock{2}{The problem}{about 55 seconds}
When a bank decides whether to grant a loan, the decision rests on one number: the probability that the applicant will default. These models are the highest-stakes statistical models in banking, because under the Basel rules and IFRS 9 they determine not just who gets credit, but how much capital the bank must hold. Now, the industry standard is remarkably old technology: a logistic-regression scorecard. Research papers keep reporting that machine learning beats it, but almost all of them measure only one thing: the area under the receiver operating characteristic curve, universally abbreviated to AUC, which is simply the ability to rank risky borrowers above safe ones. A lender needs two more things. It needs the probabilities themselves to be right, which is called calibration. And it needs the model to actually make money at its accept-or-reject cutoff. So my question is: does the machine-learning advantage survive all three tests?

\slideblock{3}{Tools and technologies}{about 30 seconds}
The whole project is Python 3.13, running on an ordinary laptop CPU. Pandas and NumPy handle the data, scikit-learn, XGBoost and LightGBM provide the models, and everything lives in a modular package on GitHub, with Jupyter, Git and LaTeX around it. One command reproduces the entire study, all seeds fixed, in about 18 minutes. Every number I will show you comes from that one command.

\slideblock{4}{Dataset}{about 75 seconds}
The data is the public Lending Club export: 621,022 personal loans issued between 2007 and 2015. Briefly, what these loans are: small unsecured consumer loans, from \$500 up to the platform's \$35,000 ceiling, median \$10,000, about \$342 a month over three years. Rates run from 5.3\% to 29\%. The typical borrower earns \$60,000 a year with a debt-to-income ratio, that is monthly debt payments divided by monthly income, of 17\%. And overwhelmingly they are refinancing: 57\% debt consolidation, another 24\% paying off credit cards. One number is worth pausing on: the FICO score, the standard American credit score produced by the Fair Isaac Corporation, running from 300 to 850. From 2009 onwards, the lowest FICO score in every single vintage is exactly 662. An identical minimum seven years running is not a coincidence, it is Lending Club's approval cutoff. So nobody genuinely subprime is in this sample at all, the median is only 692, and that truncation is part of why the achievable accuracy here sits near 0.68 rather than higher. Two design decisions matter. First, I only use 36-month loans, so every loan has finished its life before the data ends; no outcome is unknown. Second, and this was the hardest work in the project: the raw export has 151 columns, and most describe what happened after the loan was granted; payments received, recoveries, hardship flags. Any of those trivially predicts the outcome. That is called leakage, and it produces impressive, meaningless results. So I enforced a whitelist of 23 fields known on the day of application. Finally, the split follows the calendar: train on 2007 to 2013, calibrate on 2014, test on 2015. Notice that the default rate rises across those vintages. That drift is not a nuisance; it is exactly what a deployed model faces.

\slideblock{5}{Approach and code}{about 75 seconds}
Four models compete, inside one pipeline called creditrisk. The baseline is not a strawman: it is a proper industry-style scorecard. Each feature is binned, each bin gets a weight of evidence, weak features are screened out by information value, and a logistic regression on top produces a score where 20 points double the odds of repayment; the format banks actually deploy. Against it: XGBoost and LightGBM, which are gradient-boosted tree ensembles, and a small neural network. Crucially, all three challengers see exactly the same 23 features as the scorecard, through one shared design matrix, so any performance gap comes from the learner, not the preprocessing. At the bottom of the slide is the leakage decision made concrete. On the left, what I kept: loan amount, term, rate, grade, income, debt-to-income, FICO, revolving utilisation, recent inquiries; everything knowable on the day of application. On the right, what I threw away: total payments received, principal received, recoveries, last payment amount. Look at principal received to date. For every loan that was repaid, it equals the loan amount exactly. A model given that column scores almost perfectly and has learned absolutely nothing. Every model is then recalibrated the same way on the 2014 vintage, and the profit analysis uses each test loan's real cash flows: a repaid loan earns its interest, a charged-off loan loses 65\% of principal.

\slideblock{6}{Evaluation}{about 50 seconds}
Before the numbers, the metrics; three questions, each with its own. Discrimination: can the model rank bad borrowers above good ones? AUC is the probability that a random defaulter gets a higher risk score than a random good borrower, and the Kolmogorov-Smirnov statistic, KS, is the maximum distance between the two score distributions. Calibration: among loans given, say, 10\% predicted risk, do 10\% actually default? The Brier score and the expected calibration error measure that, and reliability diagrams draw it. And decision quality: realized profit per 1,000 applicants, across every possible cutoff. Throughout, I use the DeLong test and bootstrap confidence intervals to make sure any gap is real and not luck. The key point: these three questions can rank the models differently. And they do.

\slideblock{7}{Results: discrimination}{about 55 seconds}
First result: the trees win, the neural network loses. XGBoost beats the scorecard by 0.85 AUC points, LightGBM by 0.72, and with 283,026 test loans, those gaps are statistically overwhelming; the p-values are astronomically small. But look at the size: not the 8 or 15 point gaps you sometimes read from random-split studies; less than one point. And the neural network is significantly worse than the forty-year-old scorecard. Why so close? Two reasons. The features include the platform's own risk pricing, its grade and interest rate, so every model is partly re-ranking loans that were already ranked once. And the out-of-time split denies the flexible models the flattery of a random split, where training and test share the same economy.

\slideblock{8}{Results: calibration}{about 55 seconds}
Second result: calibration. In these reliability diagrams, a perfectly calibrated model lies on the diagonal. Every grey curve, the raw models, sits above it: they all under-predict risk, because the 2015 test vintage simply defaulted more than the training years. The raw trees are the worst offenders; boosted ensembles optimize ranking, not probabilities. Isotonic recalibration, fitted on 2014, pulls everyone toward the diagonal, but it cannot fix everything, because a correction learned on the past cannot anticipate drift that has not happened yet. That is precisely why regulators demand ongoing monitoring of these models. And here is the nuance I find most interesting: neither tree ensemble ever catches the scorecard on calibration, before or after recalibration. The trees keep the better Brier score, which also rewards ranking, but on pure calibration the old technology holds its ground.

\slideblock{9}{Results: profit}{about 55 seconds}
Third result: money. These curves show realized profit per 1,000 applicants as the acceptance cutoff sweeps from strict to lenient. First observation: accepting everyone already earns about \$705 per applicant, because interest income dominates this portfolio. So the model's job is only to prune the riskiest tail. The scorecard's best policy rejects about 7\% of applicants and gains 1.2\%. XGBoost rejects 9\% and gains 3.3\%. The whole battle happens in the zoomed panel on the right; everywhere else the models are interchangeable. The difference between them is about \$14,700 per 1,000 applicants, or 2.1\%. Modest per applicant, but it scales linearly with volume, which is exactly why large lenders adopt gradient boosting and small ones often, rationally, do not.

\slideblock{10}{Results: interpretability}{about 45 seconds}
What did the models actually learn? Mostly the same thing. On the left, XGBoost's feature importance; on the right, the scorecard's information values. Both agree that the platform's own pricing, the interest rate and the grade, carries the most signal. They disagree mainly in encoding: the scorecard leans on the sub-grade and FICO bins, while XGBoost prefers the continuous interest rate and income; nearly the same information, expressed differently. The rank correlation between the two is 0.56. The real difference is the interpretability trade-off: the scorecard is a readable sum of points that a loan officer can audit line by line; the ensemble needs post-hoc explanation tools.

\slideblock{11}{Ablation: hiding the platform's pricing}{about 60 seconds}
Now the follow-up experiment I think matters most. Three of my features, the grade, the sub-grade and the interest rate, are really one variable in three notations: Lending Club's own risk assessment. Within a single sub-grade, interest rates vary by only a third of a percentage point. And together they supply more than a third of the trees' predictive power. They are legitimate inputs, an investor on the platform genuinely sees them, but they narrow the question: are my models assessing credit risk, or just re-ranking loans Lending Club already ranked? So I reran everything without them. Variant A drops those three. Variant B also drops the monthly instalment, because with the term fixed at 36 months the instalment tracks the interest rate at a correlation of 0.998, so Variant A was still leaking the price through the back door. Two findings. First: every model gets worse, but the scorecard gets worse almost twice as fast as the trees, so XGBoost's lead grows from 0.83 to 2.39 AUC points. The pricing features were hiding the machine-learning advantage, not creating it. Second, and my favourite result in the project: the scorecard scores exactly the same under both variants, because its information-value screen had already thrown the instalment out as too weak. The old-fashioned screening step was what protected it from a proxy both tree ensembles were quietly feeding on.

\slideblock{12}{Robustness and limitations}{about 50 seconds}
Is this a Lending Club artifact? To check, I repeated the full protocol on an independent dataset, the UCI Taiwan credit-card data: 30,000 clients. Same ranking: XGBoost beats the scorecard by 1.50 AUC points, and the neural network again lands below it. And the honest limitations. The data contain only accepted loans, so every model inherits the platform's own selection; correcting that, reject inference, is an open research problem. The profit model is simple: undiscounted, with a fixed loss given default of 65\%. Keeping the platform's grades as features is realistic for an investor, and the ablation you just saw quantifies the understatement at roughly a factor of 3. And 2015 was a benign year; a recession vintage could reorder the calibration findings.

\slideblock{13}{Wrap-up}{about 55 seconds}
So, does machine learning beat the credit scorecard? A qualified yes. The boosted trees win on all three criteria: discrimination, Brier score after identical recalibration, and profit. The result is statistically unambiguous and it reproduces on independent data. But the margins are modest, only the trees deliver them, the neural network loses outright, neither tree ensemble ever matches the scorecard's calibration, and the size of the margin turns out to depend on the information set more than on the models. What did I learn? That leakage control was the highest-leverage work in the entire project. That the model ranking changes with the question you ask, so ``best model'' is undefined until the deployment question is fixed. And that vintage drift is why these models are monitored, not trusted. Everything you have seen, code, results, and the full report, reproduces with one command from the repository on screen. Thank you for watching.

\vspace{6mm}
\noindent\textit{Total: about 1,750 words; at 145 to 150 words per minute this is about 12 minutes. If a take runs long, the easiest cuts are the final sentences of slides 7 and 9, and the second half of slide 10.}

\end{document}
